% ============================================================================
% PRCL Intermediate Report — LaTeX
% ============================================================================
\documentclass[11pt,a4paper]{article}

% ── Packages ────────────────────────────────────────────────────────────────
\usepackage[utf8]{inputenc}
\usepackage[T1]{fontenc}
\usepackage{lmodern}
\usepackage[margin=1in]{geometry}
\usepackage{amsmath,amssymb,amsthm}
\usepackage{booktabs}
\usepackage{graphicx}
\usepackage{hyperref}
\usepackage{xcolor}
\usepackage{algorithm}
\usepackage{algpseudocode}
\usepackage{enumitem}
\usepackage{caption}
\usepackage{subcaption}
\usepackage{multirow}
\usepackage{tabularx}
\usepackage{listings}
\usepackage{fancyhdr}
\usepackage{natbib}
\usepackage{tcolorbox}

% ── Custom colors and boxes ─────────────────────────────────────────────────
\definecolor{madeup}{RGB}{180,60,60}
\definecolor{codebg}{RGB}{245,245,245}

\newtcolorbox{madeupbox}{
  colback=red!5!white,
  colframe=red!60!black,
  title={\textbf{Note: Made-Up Results}},
  fonttitle=\small,
  coltitle=white
}

\newcommand{\madeupresult}[1]{\textcolor{madeup}{#1\textsuperscript{$\dagger$}}}
\newcommand{\daggernote}{\textsuperscript{$\dagger$}\,\textcolor{madeup}{Made-up result (illustrative; pending real experiments).}}

% ── Code listing style ──────────────────────────────────────────────────────
\lstset{
  backgroundcolor=\color{codebg},
  basicstyle=\ttfamily\small,
  breaklines=true,
  frame=single,
  rulecolor=\color{gray!40},
  xleftmargin=1em,
  framexleftmargin=0.5em,
  columns=fullflexible,
  keepspaces=true,
}

% ── Header / Footer ────────────────────────────────────────────────────────
\pagestyle{fancy}
\fancyhf{}
\lhead{\small PRCL Intermediate Report}
\rhead{\small March 2026}
\cfoot{\thepage}

% ── Math operators ──────────────────────────────────────────────────────────
\DeclareMathOperator{\cosim}{cos\_sim}
\DeclareMathOperator{\softmax}{softmax}

% ============================================================================
\title{
  \textbf{PRCL: Poison-Robust Contrastive Learning\\for Vision Self-Supervised Pretraining}\\[0.5em]
  \large Intermediate Project Report
}
\author{
  Aayush Kumar\\
  \texttt{aayushakumar@github}
}
\date{March 8, 2026}

\begin{document}
\maketitle

\begin{madeupbox}
All numerical results in this report are \textbf{illustrative made-up values}
(marked with \madeupresult{red text\textsuperscript{$\dagger$}}). They represent
realistic expectations based on the system design and the SSL/backdoor defense
literature, but have \emph{not} been obtained by running the actual experiments.
Final results will replace these once GPU experiments are completed.
\end{madeupbox}

\vspace{0.5em}

% ============================================================================
\begin{abstract}
Self-supervised contrastive learning (SSL) has become the dominant paradigm for
learning visual representations from unlabeled data.  However, SSL's reliance on
large, uncurated datasets introduces a critical vulnerability: an adversary who
controls even a small fraction of the training corpus can implant backdoor
behavior that persists through downstream fine-tuning.  We present
\textbf{Poison-Robust Contrastive Learning (PRCL)}, a defense framework that
operates entirely during the pretraining phase---without requiring labels, a
clean validation set, or knowledge of the attack---to mitigate data poisoning
and backdoor attacks against contrastive SSL.

PRCL integrates three complementary mechanisms: (1)~\textbf{Probe-Consistency
Forensics (PCF)}, a label-free suspicion scoring system that identifies
anomalous samples by measuring representation stability under targeted
perturbations; (2)~\textbf{robust contrastive optimization} that reweights both
positive and negative terms in the InfoNCE objective using PCF scores, bounding
the influence of any single sample; and (3)~an optional \textbf{quarantine
buffer} that temporarily excludes highly suspicious samples and periodically
re-evaluates them as representations improve.

We evaluate PRCL on CIFAR-10, CIFAR-100, and STL-10 using SimCLR as the base
SSL method, against two families of backdoor attacks (patch-based and
blend-based) across contamination ratios from 0.5\% to 10\%.  We provide
\textbf{IntegritySuite}, a companion evaluation harness that standardizes threat
model specification, metric collection, and reproducible reporting.

This intermediate report describes the complete system design, implementation
status, and illustrative (made-up) results that set expectations for the
forthcoming full evaluation.
\end{abstract}

\textbf{Keywords:} self-supervised learning, contrastive learning, backdoor
defense, data poisoning, robust optimization

\tableofcontents
\newpage

% ============================================================================
\section{Introduction}
\label{sec:intro}

\subsection{Motivation}

The success of vision SSL methods such as SimCLR~\cite{chen2020simclr},
MoCo~\cite{he2020moco}, and BYOL~\cite{grill2020byol} has fundamentally changed
how visual representations are learned.  By pretraining on large pools of
unlabeled images, these methods produce feature extractors that rival or exceed
their supervised counterparts on downstream tasks.  However, this reliance on
unlabeled data---often scraped from the web without human
verification---introduces an underappreciated attack surface.

An adversary who can inject a small fraction~$\rho$ of poisoned samples into the
pretraining corpus can implant a \emph{backdoor}: the learned encoder behaves
normally on clean inputs, but produces representations that cause downstream
classifiers to systematically misclassify inputs containing a trigger pattern.
This attack is particularly insidious because:
\begin{enumerate}[nosep]
  \item \textbf{No labels are available during pretraining}, so label-based
    defenses (e.g., spectral signatures on label distributions) do not apply.
  \item \textbf{The defender has no clean validation set}---in the SSL setting,
    all data is assumed to be unlabeled and uncurated.
  \item \textbf{Clean utility remains high}, making the backdoor difficult to
    detect by standard evaluation protocols.
\end{enumerate}

\subsection{The Gap We Address}

Existing defenses primarily target the supervised learning setting and assume
access to labels, a clean validation set, or knowledge of the attack trigger.
In the self-supervised pretraining setting, the community lacks a
\emph{practical, label-free defense} that:
\begin{itemize}[nosep]
  \item Operates during pretraining without modifying the SSL framework fundamentally
  \item Is compatible with standard contrastive methods (SimCLR, MoCo)
  \item Preserves clean downstream accuracy (utility)
  \item Provides interpretable forensic signals
  \item Ships with reproducible evaluation infrastructure
\end{itemize}

\subsection{Contributions}

We make the following contributions:

\begin{enumerate}
  \item \textbf{PRCL (Poison-Robust Contrastive Learning):} A plug-and-play
    defense framework for contrastive SSL pretraining that combines probe-based
    forensics with robust contrastive optimization.  PRCL is fully label-free
    and does not require a separate clean validation set.

  \item \textbf{Probe-Consistency Forensics (PCF):} A novel per-sample suspicion
    scoring mechanism that exploits the observation that backdoor-poisoned
    samples exhibit abnormal representation instability under targeted probe
    perturbations (blur, occlusion, frequency filtering, desaturation).

  \item \textbf{Robust Weighted InfoNCE:} A modified contrastive objective that
    uses PCF suspicion scores to reweight both positive pair contributions and
    negative sample influence, bounding the impact of any individual sample on
    the learned representation.

  \item \textbf{IntegritySuite:} An evaluation harness for SSL backdoor defense
    research that standardizes threat model specification, attack/defense
    configuration, metric collection, and reproducible reporting.

  \item \textbf{Empirical evaluation} on CIFAR-10, CIFAR-100, and STL-10 against
    patch-based and blend-based backdoor attacks at contamination ratios
    $\rho\in\{0.5\%,1\%,3\%,5\%,10\%\}$.
\end{enumerate}


% ============================================================================
\section{Related Work}
\label{sec:related}

\subsection{Self-Supervised Contrastive Learning}

Contrastive SSL learns visual representations by maximizing agreement between
differently augmented views of the same image while pushing apart
representations of different images.  \textbf{SimCLR}~\cite{chen2020simclr}
introduced the NT-Xent loss over augmented pairs within a mini-batch.
\textbf{MoCo}~\cite{he2020moco} extended this with a momentum encoder and
queue-based negative storage.  Subsequent work (BYOL, Barlow Twins, VICReg)
explored non-contrastive objectives, but the contrastive framework remains the
most widely deployed.

The InfoNCE loss~\cite{oord2018cpc}, which underlies both SimCLR and MoCo, for a
batch of $N$ samples with two views each is:
\begin{equation}
  \mathcal{L}_{\text{InfoNCE}}
  = -\frac{1}{2N}\sum_{i=1}^{2N}
    \log\frac{\exp\bigl(\text{sim}(z_i, z_{p(i)})/\tau\bigr)}
             {\sum_{j\neq i}\exp\bigl(\text{sim}(z_i, z_j)/\tau\bigr)}
  \label{eq:infonce}
\end{equation}
where $z_i$ are $\ell_2$-normalized projected features, $p(i)$ is the positive
pair index, $\text{sim}(\cdot,\cdot)$ is cosine similarity, and $\tau$ is the
temperature parameter.

\subsection{Backdoor Attacks on SSL}

Backdoor attacks against SSL have been demonstrated by several recent works.
\textbf{SSL-Backdoor}~\cite{saha2022backdoor} showed that poisoning a small
fraction of the pretraining data with trigger-stamped images can implant
persistent backdoor behavior.  \textbf{BadEncoder}~\cite{jia2022badencoder}
extended this to contrastive encoders specifically.  Two primary trigger families
are well-established:

\begin{itemize}[nosep]
  \item \textbf{Patch-based (BadNets-style)~\cite{gu2017badnets}:} A small
    solid-color patch placed at a fixed image position.  Creates a localized,
    high-frequency artifact.
  \item \textbf{Blend-based:} A global trigger pattern alpha-blended across the
    entire image.  Creates a distributed, lower-frequency artifact that is
    harder to detect visually.
\end{itemize}

\subsection{Backdoor Defenses}

Most defenses target supervised learning and assume access to labels:
\textbf{Spectral Signatures}~\cite{tran2018spectral} detects poisoned samples
via top singular vectors;
\textbf{Neural Cleanse}~\cite{wang2019neural} reverse-engineers triggers;
\textbf{ABL}~\cite{li2021anti} uses loss-guided sample selection.
\textbf{DECREE}~\cite{zhang2023decree} is among the few targeting SSL
specifically.

\textbf{Gap:} These methods either require labels, assume a clean validation
set, or operate post-training.  PRCL operates during self-supervised pretraining
itself, using only the unlabeled training data.

\subsection{Robust Estimation}

PRCL draws on classical ideas from robust statistics, particularly
\emph{bounded-influence estimators}~\cite{hampel1986robust} and \emph{trimmed
means}~\cite{tukey1960survey}.  In a contamination model where
$X\sim(1-\rho)P+\rho Q$, robust estimators control the maximum influence of any
single observation.  PRCL applies this principle to the InfoNCE gradient by
reweighting sample contributions based on suspicion scores.


% ============================================================================
\section{Threat Model}
\label{sec:threat}

\subsection{Defender Model}

The defender controls the SSL pretraining pipeline:
\begin{itemize}[nosep]
  \item Choice of augmentations, sampling strategy, loss function, optimizer
  \item Access to training data but \textbf{no labels} during pretraining
  \item Ability to modify training dynamics (sample weighting, gradient clipping)
  \item Does \textbf{not} know which specific samples are poisoned
  \item Does \textbf{not} have a separate clean validation set
  \item Does \textbf{not} know the trigger pattern, target class, or attack family
\end{itemize}

\subsection{Attacker Model}

The attacker can inject a fraction $\rho$ of poisoned samples:
\begin{itemize}[nosep]
  \item \textbf{Capability:} Modify
    $\rho\in\{0.005,0.01,0.03,0.05,0.10\}$ of the training data
  \item \textbf{Trigger:} Apply a function $T:\mathcal{X}\to\mathcal{X}$ to
    each poisoned sample
  \item \textbf{Objective:} Downstream classifiers misclassify triggered inputs
    as target class~$y_t$ while maintaining high clean accuracy
\end{itemize}

\subsection{Evaluation Tiers}

\begin{table}[h]
\centering
\caption{Attacker knowledge tiers.}
\begin{tabular}{@{}lll@{}}
\toprule
\textbf{Tier} & \textbf{Knowledge} & \textbf{Description} \\
\midrule
A (Gray-box)  & SSL method, not defense & Standard setting \\
B (White-box) & Full training setup      & Defense-aware tuning \\
C (Adaptive)  & PRCL internals           & Discussion only \\
\bottomrule
\end{tabular}
\label{tab:tiers}
\end{table}

Our primary evaluation is at Tier~A (gray-box), with a discussion of adaptive
attacks (Tier~C) in Section~\ref{sec:limitations}.

\subsection{Attack Families}

\begin{table}[h]
\centering
\caption{Attack families evaluated.}
\begin{tabular}{@{}lll@{}}
\toprule
\textbf{Attack} & \textbf{Mechanism} & \textbf{Trigger Type} \\
\midrule
Patch Backdoor & $4\!\times\!4$ white patch, bottom-right & Localized, high-freq \\
Blend Backdoor & Full-image pattern, $\alpha=0.1$         & Global, distributed \\
\bottomrule
\end{tabular}
\label{tab:attacks}
\end{table}


% ============================================================================
\section{Method: Poison-Robust Contrastive Learning}
\label{sec:method}

PRCL integrates three complementary defense mechanisms into the standard
contrastive SSL training loop.  The key insight is that \emph{poisoned samples
exhibit abnormal representation behavior when subjected to targeted
perturbations}---perturbations designed to disrupt trigger artifacts while
preserving natural image semantics.

\subsection{Overview}

Given a training corpus $\mathcal{D}=\{x_1,\dots,x_N\}$ containing a fraction
$\rho$ of poisoned samples, PRCL modifies the SimCLR training pipeline as
follows:
\begin{enumerate}[nosep]
  \item \textbf{PCF:} Compute per-sample suspicion $q(x_i)\in[0,1]$ by
    measuring representation stability under probe perturbations.
  \item \textbf{Robust Loss:} Replace InfoNCE with a weighted variant using
    $q(x_i)$ to reduce the influence of suspicious samples.
  \item \textbf{Quarantine (optional):} Maintain a buffer of highly suspicious
    samples, periodically re-evaluating them.
\end{enumerate}

\subsection{PRCL Training Algorithm}

\begin{algorithm}[H]
\caption{PRCL Training}
\label{alg:prcl}
\begin{algorithmic}[1]
\Require Dataset $\mathcal{D}$, encoder~$f_\theta$, $K$ probes,
  $\lambda_{\text{pos}},\lambda_{\text{neg}},w_{\min},C$
\Ensure Trained encoder $f_\theta$
\For{epoch $=1$ to $E$}
  \For{mini-batch $B=\{(v_{1,i},v_{2,i})\}_{i=1}^{|B|}$}
    \State $h_i\gets f_\theta(v_{1,i})$; \quad $z_i\gets g(h_i)$
      \Comment{Backbone + projection}
    \For{$k=1$ to $K$}
      \Comment{PCF scoring}
      \State $h'_{i,k}\gets f_\theta(\text{Probe}_k(v_{1,i}))$
      \State $s_{i,k}\gets\cos(h_i, h'_{i,k})$
    \EndFor
    \State $q_{\text{raw}}(x_i)\gets 1-\tfrac{1}{K}\sum_k s_{i,k}$
    \State $q(x_i)\gets\text{RollingZScore}(q_{\text{raw}}(x_i))$
      \Comment{Normalize to $[0,1]$}
    \State $w_{\text{pos}}(i)\gets\max\bigl(w_{\min},\;
      1-\lambda_{\text{pos}}\cdot q(x_i)\bigr)$
    \State $w_{\text{neg}}(i,j)\gets\max\bigl(w_{\min},\;
      1-\lambda_{\text{neg}}\cdot q(x_j)\bigr)$
    \State $\mathcal{L}\gets\text{WeightedInfoNCE}(z,w_{\text{pos}},
      w_{\text{neg}},\tau)$
    \State $\mathcal{L}_{\text{capped}}\gets
      \min\bigl(w_{\text{pos}}\cdot\mathcal{L},\;C\bigr)$
      \Comment{Gradient cap}
    \State Update $\theta$ via $\nabla\mathcal{L}_{\text{capped}}$
  \EndFor
\EndFor
\end{algorithmic}
\end{algorithm}

\subsection{Probe-Consistency Forensics (PCF)}

PCF computes a per-sample suspicion score by measuring how much a sample's
learned representation changes under targeted perturbations.

\subsubsection{Intuition}

Backdoor triggers are designed to be features that the encoder learns to
associate with a target class.  When we apply perturbations that disrupt the
trigger pattern, the representation of a poisoned sample changes significantly
because the learned ``trigger feature'' is disrupted.  Clean samples, whose
representations are based on genuine semantic content, are more stable.

\subsubsection{Probe Transforms}

PRCL provides a registry of probe transforms:

\begin{table}[h]
\centering
\caption{Probe transforms for PCF forensics.}
\begin{tabular}{@{}lll@{}}
\toprule
\textbf{Probe} & \textbf{Mechanism} & \textbf{Target Artifact} \\
\midrule
Blur         & Gaussian blur ($\sigma\!=\!2.0$, $7\!\times\!7$) &
  High-freq patches \\
Occlusion    & Random 25\%-area zeroing &
  Localized patterns \\
Freq Lowpass & FFT low-pass (cutoff $=0.3$) &
  Spectral artifacts \\
Desaturation & ITU-R 601-2 grayscale &
  Color-based triggers \\
\bottomrule
\end{tabular}
\label{tab:probes}
\end{table}

\subsubsection{Probe-Alignment Scoring}

For each sample $x_i$, compute the cosine similarity between the clean embedding
and each probe embedding:
\begin{equation}
  q_{\text{raw}}(x_i) = 1 - \frac{1}{K}\sum_{k=1}^{K}
    \cos\!\bigl(h_{\text{clean}}^{(i)},\; h_{\text{probe}_k}^{(i)}\bigr)
  \label{eq:pcf}
\end{equation}
where $h_{\text{clean}}^{(i)}=f_\theta(x_i)$ and
$h_{\text{probe}_k}^{(i)}=f_\theta(\text{Probe}_k(x_i))$, both
$\ell_2$-normalized.

\subsubsection{Score Normalization}

Raw scores are normalized to $[0,1]$ using a rolling z-score normalizer with
exponential moving average (momentum $\beta=0.1$):
\begin{align}
  \hat{\mu}_t &= (1-\beta)\,\hat{\mu}_{t-1} + \beta\cdot\text{mean}(q_{\text{raw}})
  \label{eq:ema_mu} \\
  \hat{\sigma}_t^2 &= (1-\beta)\,\hat{\sigma}_{t-1}^2 + \beta\cdot\text{var}(q_{\text{raw}})
  \label{eq:ema_var} \\
  q(x_i) &= \sigma\!\left(\frac{q_{\text{raw}}(x_i)-\hat{\mu}_t}{\hat{\sigma}_t}\right)
  \label{eq:sigmoid_norm}
\end{align}

\subsection{Robust Weighted InfoNCE Loss}

PRCL replaces the standard InfoNCE loss with a weighted variant.  \textbf{All
weights are detached}---no gradient flows through suspicion scores.

\subsubsection{Positive Reweighting}
\begin{equation}
  w_{\text{pos}}(i)=\max\!\bigl(w_{\min},\;1-\lambda_{\text{pos}}\cdot q(x_i)\bigr)
  \label{eq:wpos}
\end{equation}

\subsubsection{Negative Reweighting}
\begin{equation}
  w_{\text{neg}}(i,j)=\max\!\bigl(w_{\min},\;1-\lambda_{\text{neg}}\cdot q(x_j)\bigr)
  \label{eq:wneg}
\end{equation}
applied in the log domain: $\tilde{s}_{ij}=s_{ij}+\log w_{\text{neg}}(i,j)$.

\subsubsection{Gradient Capping}
\begin{equation}
  \mathcal{L}_i^{\text{capped}}=\min\!\bigl(w_{\text{pos}}(i)\cdot\mathcal{L}_i,\;C\bigr)
  \label{eq:gradcap}
\end{equation}

\subsubsection{Combined Loss}
\begin{equation}
  \mathcal{L}_{\text{PRCL}}
  = \frac{1}{\sum_i w_{\text{pos}}(i)}
    \sum_{i=1}^{2N}\mathcal{L}_i^{\text{capped}}
  \label{eq:prcl_loss}
\end{equation}

\subsection{Quarantine Buffer (Optional)}

The quarantine mechanism is \textbf{disabled by default} and intended for
high-contamination scenarios:
\begin{enumerate}[nosep]
  \item \textbf{EMA tracking:}
    $\bar{q}_i^{(t)}=(1-\alpha)\bar{q}_i^{(t-1)}+\alpha\,q_i^{(t)}$,
    $\alpha=0.3$
  \item \textbf{Quarantine:} if $\bar{q}_i>\tau_q=0.95$ or exceeds the 99th
    percentile
  \item \textbf{Re-evaluation:} every $R=5$ epochs; reinstate if EMA drops
    below threshold
\end{enumerate}


% ============================================================================
\section{Implementation Status}
\label{sec:implementation}

The complete PRCL + IntegritySuite codebase has been implemented, tested, and
committed.  Table~\ref{tab:implementation} summarizes the implementation status.

\begin{table}[h]
\centering
\caption{Implementation status of all system components.}
\begin{tabular}{@{}llc@{}}
\toprule
\textbf{Component} & \textbf{Location} & \textbf{Status} \\
\midrule
Project bootstrap (Hydra configs, schema) & \texttt{configs/}, \texttt{config\_schema.py} & \checkmark \\
SimCLR SSL core (backbone, head, loss, loop) & \texttt{src/prcl/ssl/} & \checkmark \\
Dataset builders (CIFAR-10/100, STL-10, poison wrapper) & \texttt{src/prcl/datasets/} & \checkmark \\
PRCL defense (PCF, robust loss, quarantine) & \texttt{src/prcl/defenses/prcl/} & \checkmark \\
Attack adapters (patch, blend, safety gate) & \texttt{src/prcl/attacks/} & \checkmark \\
Evaluation (linear probe, ASR, forensics) & \texttt{src/prcl/eval/} & \checkmark \\
IntegritySuite (run manager, cards, plots, aggregate) & \texttt{src/prcl/integritysuite/} & \checkmark \\
Training entry point & \texttt{scripts/train.py} & \checkmark \\
Sweep \& ablation scripts & \texttt{scripts/*.sh} & \checkmark \\
Test suite (99 tests) & \texttt{tests/} & \checkmark \\
Documentation & \texttt{docs/} & \checkmark \\
Pre-commit \& lint (ruff, 0 violations) & \texttt{.pre-commit-config.yaml} & \checkmark \\
Colab notebook & \texttt{notebooks/colab\_train.ipynb} & \checkmark \\
\midrule
GPU experiment runs & \texttt{runs/} & Pending \\
\bottomrule
\end{tabular}
\label{tab:implementation}
\end{table}

\subsection{Test Suite}

The codebase passes \textbf{99 tests} across unit, integration, and smoke test
categories:

\begin{table}[h]
\centering
\caption{Test suite breakdown.}
\begin{tabular}{@{}lr@{}}
\toprule
\textbf{Category} & \textbf{Tests} \\
\midrule
SSL core (backbone, head, loss, transforms, train loop) & 15 \\
PRCL defense (PCF, robust loss, probes, quarantine, thresholds) & 33 \\
Attack adapters (patch, blend, builder, safety) & 17 \\
Datasets (CIFAR, STL-10, poison wrapper) & 5 \\
IntegritySuite (cards, schemas, visualizations) & 12 \\
Run manager & 7 \\
Transforms & 4 \\
Integration (end-to-end pipeline) & 3 \\
Smoke (config loading, imports, schema) & 7 \\
\midrule
\textbf{Total} & \textbf{99} \\
\bottomrule
\end{tabular}
\end{table}

\subsection{Code Quality}

\begin{itemize}[nosep]
  \item \textbf{Lint:} Ruff---zero violations across all source and test files
  \item \textbf{Pre-commit hooks:} ruff (lint + format), mypy, trailing whitespace,
    end-of-file, YAML/JSON checks
  \item \textbf{Package:} pip-installable via \texttt{pip install -e '.[dev]'}
\end{itemize}


% ============================================================================
\section{Experimental Setup}
\label{sec:setup}

\subsection{Datasets}

\begin{table}[h]
\centering
\caption{Datasets used for evaluation.}
\begin{tabular}{@{}lcccc@{}}
\toprule
\textbf{Dataset} & \textbf{Classes} & \textbf{Train Size} & \textbf{Image Size} & \textbf{SSL Use} \\
\midrule
CIFAR-10  & 10  & 50{,}000  & $32\!\times\!32$  & Primary benchmark \\
CIFAR-100 & 100 & 50{,}000  & $32\!\times\!32$  & Fine-grained \\
STL-10    & 10  & $\sim$105{,}000 & $96\!\times\!96$ & Realistic SSL \\
\bottomrule
\end{tabular}
\end{table}

\subsection{Training Configuration}

\begin{table}[h]
\centering
\caption{SimCLR training hyperparameters.}
\begin{tabular}{@{}ll@{}}
\toprule
\textbf{Parameter} & \textbf{Value} \\
\midrule
Backbone      & ResNet-18 (512-d features) \\
Projection    & MLP $512\to2048\to128$ \\
Optimizer     & Adam, $\text{lr}=3\!\times\!10^{-4}$, wd=$10^{-4}$ \\
Batch size    & 256 \\
Epochs        & 200 \\
Temperature   & $\tau=0.5$ \\
Warmup        & 10 epochs \\
LR schedule   & Cosine decay after warmup \\
\bottomrule
\end{tabular}
\end{table}

\subsection{PRCL Defense Defaults}

\begin{table}[h]
\centering
\caption{Default PRCL hyperparameters.}
\begin{tabular}{@{}lll@{}}
\toprule
\textbf{Parameter} & \textbf{Default} & \textbf{Description} \\
\midrule
PCF probes & blur, occlusion & Probe types \\
PCF statistic & probe\_alignment & Scoring strategy \\
Normalization & rolling\_zscore & Score calibration \\
$\lambda_{\text{pos}}$ & 0.5 & Positive reweighting \\
$\lambda_{\text{neg}}$ & 0.4 & Negative reweighting \\
$w_{\min}$ & 0.2 & Minimum weight floor \\
Gradient cap $C$ & 5.0 & Per-sample loss cap \\
Quarantine & disabled & Optional aggressive mode \\
\bottomrule
\end{tabular}
\end{table}

\subsection{Evaluation Protocol}

\textbf{Linear probe:} Freeze encoder; train a linear classifier (SGD,
$\text{lr}=0.1$, cosine schedule, 100 epochs).  Report top-1 test accuracy
(\textbf{ACC$\uparrow$}).

\textbf{Attack success rate:} For each non-target test image, apply trigger,
classify via encoder + linear head.  Report fraction classified as target
(\textbf{ASR$\downarrow$}).

\subsection{Metrics}

\begin{table}[h]
\centering
\caption{Evaluation metrics.}
\begin{tabular}{@{}lll@{}}
\toprule
\textbf{Metric} & \textbf{Notation} & \textbf{Description} \\
\midrule
Clean Accuracy & ACC$\uparrow$ & Linear probe test accuracy \\
Attack Success Rate & ASR$\downarrow$ & Triggered $\to$ target fraction \\
Accuracy Drop & $\Delta$ACC & ACC(defended) $-$ ACC(baseline) \\
ASR Reduction & $\Delta$ASR & ASR(undefended) $-$ ASR(defended) \\
PCF AUC & AUC$\uparrow$ & ROC-AUC of $q(x)$ vs.\ true labels \\
\bottomrule
\end{tabular}
\end{table}


% ============================================================================
\section{Results (Made-Up)}
\label{sec:results}

\begin{madeupbox}
\textbf{All results in this section are illustrative made-up values.}  They are
designed to represent realistic expectations based on the SSL/backdoor defense
literature and our system design.  Each made-up value is marked with
\madeupresult{\textsuperscript{$\dagger$}}.  Real results will be obtained by
running the experiments on GPU hardware.
\end{madeupbox}

% ── 7.1 CIFAR-10 ────────────────────────────────────────────────────────────
\subsection{Main Results: CIFAR-10}

\begin{table}[h]
\centering
\caption{CIFAR-10, ResNet-18, \textbf{Patch Backdoor} attack. \daggernote}
\label{tab:cifar10_patch}
\resizebox{\textwidth}{!}{%
\begin{tabular}{@{}l cc cc cc cc@{}}
\toprule
& \multicolumn{2}{c}{$\rho=0.5\%$}
& \multicolumn{2}{c}{$\rho=1\%$}
& \multicolumn{2}{c}{$\rho=5\%$}
& \multicolumn{2}{c}{$\rho=10\%$} \\
\cmidrule(lr){2-3}\cmidrule(lr){4-5}\cmidrule(lr){6-7}\cmidrule(lr){8-9}
\textbf{Config}
& ACC$\uparrow$ & ASR$\downarrow$
& ACC$\uparrow$ & ASR$\downarrow$
& ACC$\uparrow$ & ASR$\downarrow$
& ACC$\uparrow$ & ASR$\downarrow$ \\
\midrule
Clean baseline
  & \madeupresult{88.4} & ---
  & \madeupresult{88.4} & ---
  & \madeupresult{88.4} & ---
  & \madeupresult{88.4} & --- \\
Poisoned (no defense)
  & \madeupresult{88.1} & \madeupresult{34.2}
  & \madeupresult{87.9} & \madeupresult{58.7}
  & \madeupresult{87.2} & \madeupresult{82.4}
  & \madeupresult{86.5} & \madeupresult{91.3} \\
\textbf{PRCL (ours)}
  & \madeupresult{87.8} & \madeupresult{12.1}
  & \madeupresult{87.5} & \madeupresult{18.3}
  & \madeupresult{86.4} & \madeupresult{35.6}
  & \madeupresult{85.1} & \madeupresult{52.8} \\
\bottomrule
\end{tabular}%
}
\end{table}

\begin{table}[h]
\centering
\caption{CIFAR-10, ResNet-18, \textbf{Blend Backdoor} attack. \daggernote}
\label{tab:cifar10_blend}
\resizebox{\textwidth}{!}{%
\begin{tabular}{@{}l cc cc cc cc@{}}
\toprule
& \multicolumn{2}{c}{$\rho=0.5\%$}
& \multicolumn{2}{c}{$\rho=1\%$}
& \multicolumn{2}{c}{$\rho=5\%$}
& \multicolumn{2}{c}{$\rho=10\%$} \\
\cmidrule(lr){2-3}\cmidrule(lr){4-5}\cmidrule(lr){6-7}\cmidrule(lr){8-9}
\textbf{Config}
& ACC$\uparrow$ & ASR$\downarrow$
& ACC$\uparrow$ & ASR$\downarrow$
& ACC$\uparrow$ & ASR$\downarrow$
& ACC$\uparrow$ & ASR$\downarrow$ \\
\midrule
Clean baseline
  & \madeupresult{88.4} & ---
  & \madeupresult{88.4} & ---
  & \madeupresult{88.4} & ---
  & \madeupresult{88.4} & --- \\
Poisoned (no defense)
  & \madeupresult{88.2} & \madeupresult{28.5}
  & \madeupresult{88.0} & \madeupresult{47.1}
  & \madeupresult{87.5} & \madeupresult{73.8}
  & \madeupresult{86.9} & \madeupresult{85.2} \\
\textbf{PRCL (ours)}
  & \madeupresult{87.9} & \madeupresult{15.7}
  & \madeupresult{87.4} & \madeupresult{23.9}
  & \madeupresult{86.7} & \madeupresult{42.3}
  & \madeupresult{85.6} & \madeupresult{58.1} \\
\bottomrule
\end{tabular}%
}
\end{table}

\textbf{Observations (expected):}
\begin{itemize}[nosep]
  \item PRCL reduces ASR substantially at low-to-moderate contamination
    ($\rho\leq 5\%$) while maintaining clean accuracy within $\sim$2\% of
    baseline.
  \item Patch-based attacks show higher ASR ceilings than blend-based at the
    same $\rho$ (patch triggers are stronger), but PRCL also achieves larger
    ASR reductions against them (because blur/occlusion probes are effective at
    disrupting localized patches).
  \item At $\rho=10\%$, defense effectiveness degrades, consistent with the
    bounded-influence limitation.
\end{itemize}


% ── 7.2 CIFAR-100 ───────────────────────────────────────────────────────────
\subsection{Main Results: CIFAR-100}

\begin{table}[h]
\centering
\caption{CIFAR-100, ResNet-18, \textbf{Patch Backdoor} attack. \daggernote}
\label{tab:cifar100_patch}
\begin{tabular}{@{}l cc cc cc@{}}
\toprule
& \multicolumn{2}{c}{$\rho=1\%$}
& \multicolumn{2}{c}{$\rho=5\%$}
& \multicolumn{2}{c}{$\rho=10\%$} \\
\cmidrule(lr){2-3}\cmidrule(lr){4-5}\cmidrule(lr){6-7}
\textbf{Config}
& ACC$\uparrow$ & ASR$\downarrow$
& ACC$\uparrow$ & ASR$\downarrow$
& ACC$\uparrow$ & ASR$\downarrow$ \\
\midrule
Clean baseline
  & \madeupresult{61.2} & ---
  & \madeupresult{61.2} & ---
  & \madeupresult{61.2} & --- \\
Poisoned (no defense)
  & \madeupresult{60.8} & \madeupresult{52.3}
  & \madeupresult{60.1} & \madeupresult{78.9}
  & \madeupresult{58.7} & \madeupresult{89.1} \\
\textbf{PRCL (ours)}
  & \madeupresult{60.2} & \madeupresult{16.4}
  & \madeupresult{59.0} & \madeupresult{33.7}
  & \madeupresult{57.4} & \madeupresult{51.2} \\
\bottomrule
\end{tabular}
\end{table}

\begin{table}[h]
\centering
\caption{CIFAR-100, ResNet-18, \textbf{Blend Backdoor} attack. \daggernote}
\label{tab:cifar100_blend}
\begin{tabular}{@{}l cc cc cc@{}}
\toprule
& \multicolumn{2}{c}{$\rho=1\%$}
& \multicolumn{2}{c}{$\rho=5\%$}
& \multicolumn{2}{c}{$\rho=10\%$} \\
\cmidrule(lr){2-3}\cmidrule(lr){4-5}\cmidrule(lr){6-7}
\textbf{Config}
& ACC$\uparrow$ & ASR$\downarrow$
& ACC$\uparrow$ & ASR$\downarrow$
& ACC$\uparrow$ & ASR$\downarrow$ \\
\midrule
Clean baseline
  & \madeupresult{61.2} & ---
  & \madeupresult{61.2} & ---
  & \madeupresult{61.2} & --- \\
Poisoned (no defense)
  & \madeupresult{61.0} & \madeupresult{41.6}
  & \madeupresult{60.4} & \madeupresult{68.2}
  & \madeupresult{59.5} & \madeupresult{80.7} \\
\textbf{PRCL (ours)}
  & \madeupresult{60.4} & \madeupresult{19.8}
  & \madeupresult{59.3} & \madeupresult{38.5}
  & \madeupresult{57.9} & \madeupresult{55.4} \\
\bottomrule
\end{tabular}
\end{table}


% ── 7.3 STL-10 ──────────────────────────────────────────────────────────────
\subsection{Main Results: STL-10}

\begin{table}[h]
\centering
\caption{STL-10, ResNet-18, \textbf{Patch Backdoor} attack. \daggernote}
\label{tab:stl10}
\begin{tabular}{@{}l cc cc@{}}
\toprule
& \multicolumn{2}{c}{$\rho=1\%$}
& \multicolumn{2}{c}{$\rho=5\%$} \\
\cmidrule(lr){2-3}\cmidrule(lr){4-5}
\textbf{Config}
& ACC$\uparrow$ & ASR$\downarrow$
& ACC$\uparrow$ & ASR$\downarrow$ \\
\midrule
Clean baseline
  & \madeupresult{82.6} & ---
  & \madeupresult{82.6} & --- \\
Poisoned (no defense)
  & \madeupresult{82.1} & \madeupresult{55.0}
  & \madeupresult{81.0} & \madeupresult{79.3} \\
\textbf{PRCL (ours)}
  & \madeupresult{81.4} & \madeupresult{19.2}
  & \madeupresult{80.1} & \madeupresult{37.8} \\
\bottomrule
\end{tabular}
\end{table}


% ── 7.4 PCF Quality ─────────────────────────────────────────────────────────
\subsection{PCF Forensic Quality}

\begin{table}[h]
\centering
\caption{PCF detection quality: ROC-AUC of suspicion scores vs.\ true poison labels. \daggernote}
\label{tab:pcf_auc}
\begin{tabular}{@{}llccc@{}}
\toprule
\textbf{Attack} & \textbf{Dataset}
& $\rho=1\%$ & $\rho=5\%$ & $\rho=10\%$ \\
\midrule
Patch Backdoor & CIFAR-10
  & \madeupresult{0.91} & \madeupresult{0.88} & \madeupresult{0.83} \\
Patch Backdoor & CIFAR-100
  & \madeupresult{0.89} & \madeupresult{0.85} & \madeupresult{0.80} \\
Blend Backdoor & CIFAR-10
  & \madeupresult{0.82} & \madeupresult{0.78} & \madeupresult{0.73} \\
Blend Backdoor & CIFAR-100
  & \madeupresult{0.79} & \madeupresult{0.74} & \madeupresult{0.69} \\
\bottomrule
\end{tabular}
\end{table}

\textbf{Observations (expected):}
\begin{itemize}[nosep]
  \item PCF achieves higher AUC for patch-based attacks ($\sim$0.88--0.91)
    because blur and occlusion probes are particularly effective at disrupting
    localized high-frequency patches.
  \item Blend attacks are harder to detect ($\sim$0.73--0.82 AUC) because the
    distributed trigger is more resilient to probe perturbations.
  \item AUC decreases at higher $\rho$ as the larger poisoned fraction shifts
    the score distribution.
\end{itemize}


% ── 7.5 Utility ─────────────────────────────────────────────────────────────
\subsection{Utility Preservation}

\begin{table}[h]
\centering
\caption{Clean accuracy comparison (no attack). \daggernote}
\label{tab:utility}
\begin{tabular}{@{}llccc@{}}
\toprule
\textbf{Dataset} & \textbf{Backbone}
& \textbf{No Defense} & \textbf{PRCL} & $\Delta$\textbf{ACC} \\
\midrule
CIFAR-10  & ResNet-18
  & \madeupresult{88.4} & \madeupresult{87.1} & \madeupresult{$-1.3$} \\
CIFAR-100 & ResNet-18
  & \madeupresult{61.2} & \madeupresult{59.8} & \madeupresult{$-1.4$} \\
STL-10    & ResNet-18
  & \madeupresult{82.6} & \madeupresult{81.2} & \madeupresult{$-1.4$} \\
\bottomrule
\end{tabular}
\end{table}

PRCL's overhead on clean accuracy is expected to be $\sim$1--2\%, which is
acceptable for security-critical applications.


% ============================================================================
\section{Ablation Studies (Made-Up)}
\label{sec:ablation}

\begin{madeupbox}
All ablation results below are illustrative made-up values.
\end{madeupbox}

\subsection{Component Ablation}

\begin{table}[h]
\centering
\caption{Component ablation on CIFAR-10, Patch Backdoor, $\rho=1\%$, ResNet-18. \daggernote}
\label{tab:ablation}
\begin{tabular}{@{}llcc@{}}
\toprule
\textbf{Configuration} & \textbf{Components}
& ACC$\uparrow$ & ASR$\downarrow$ \\
\midrule
No defense & ---
  & \madeupresult{87.9} & \madeupresult{58.7} \\
PCF only & PCF scoring (no reweighting)
  & \madeupresult{87.7} & \madeupresult{51.2} \\
PCF + pos.\ reweight & PCF + $w_{\text{pos}}$
  & \madeupresult{87.6} & \madeupresult{34.8} \\
PCF + neg.\ sanitize & PCF + $w_{\text{neg}}$
  & \madeupresult{87.5} & \madeupresult{39.1} \\
PRCL w/o grad cap & PCF + $w_{\text{pos}}$ + $w_{\text{neg}}$
  & \madeupresult{87.4} & \madeupresult{22.6} \\
\textbf{PRCL (full)} & PCF + $w_{\text{pos}}$ + $w_{\text{neg}}$ + grad cap
  & \madeupresult{87.5} & \madeupresult{18.3} \\
\bottomrule
\end{tabular}
\end{table}

\textbf{Observations (expected):}
\begin{itemize}[nosep]
  \item PCF scoring alone (without reweighting) already provides mild ASR
    reduction, suggesting that the scoring signal is informative.
  \item Positive reweighting is the single most impactful component
    ($\sim$24pp ASR reduction).
  \item Negative sanitization provides complementary benefit ($\sim$5pp
    additional reduction when combined with positive reweighting).
  \item Gradient capping provides a modest further improvement ($\sim$4pp)
    and acts as a safety net for outlier batches.
\end{itemize}

\subsection{Sensitivity Analysis: $\lambda_{\text{pos}}$}

\begin{table}[h]
\centering
\caption{$\lambda_{\text{pos}}$ sensitivity (CIFAR-10, Patch Backdoor, $\rho\!=\!1\%$). \daggernote}
\label{tab:lambda_pos}
\begin{tabular}{@{}ccc@{}}
\toprule
$\lambda_{\text{pos}}$ & ACC$\uparrow$ & ASR$\downarrow$ \\
\midrule
0.0 & \madeupresult{87.9} & \madeupresult{52.4} \\
0.2 & \madeupresult{87.8} & \madeupresult{32.1} \\
\textbf{0.5 (default)} & \madeupresult{87.5} & \madeupresult{18.3} \\
0.8 & \madeupresult{86.8} & \madeupresult{14.7} \\
1.0 & \madeupresult{85.9} & \madeupresult{13.2} \\
\bottomrule
\end{tabular}
\end{table}

Higher $\lambda_{\text{pos}}$ yields lower ASR but at increasing clean accuracy
cost.  The default $\lambda_{\text{pos}}=0.5$ balances defense and utility.

\subsection{Sensitivity Analysis: $\lambda_{\text{neg}}$}

\begin{table}[h]
\centering
\caption{$\lambda_{\text{neg}}$ sensitivity (CIFAR-10, Patch Backdoor, $\rho\!=\!1\%$). \daggernote}
\label{tab:lambda_neg}
\begin{tabular}{@{}ccc@{}}
\toprule
$\lambda_{\text{neg}}$ & ACC$\uparrow$ & ASR$\downarrow$ \\
\midrule
0.0 & \madeupresult{87.6} & \madeupresult{25.4} \\
0.2 & \madeupresult{87.5} & \madeupresult{21.8} \\
\textbf{0.4 (default)} & \madeupresult{87.5} & \madeupresult{18.3} \\
0.6 & \madeupresult{87.2} & \madeupresult{16.9} \\
0.8 & \madeupresult{86.7} & \madeupresult{16.1} \\
\bottomrule
\end{tabular}
\end{table}

Negative reweighting shows diminishing returns beyond $\lambda_{\text{neg}}=0.4$
while increasingly affecting clean accuracy.

\subsection{Multi-Seed Variance}

\begin{table}[h]
\centering
\caption{Multi-seed results (CIFAR-10, Patch Backdoor, $\rho\!=\!1\%$, 3 seeds). \daggernote}
\label{tab:multiseed}
\begin{tabular}{@{}lcccc@{}}
\toprule
\textbf{Metric} & \textbf{Seed 42} & \textbf{Seed 123} & \textbf{Seed 456}
& \textbf{Mean $\pm$ Std} \\
\midrule
Clean ACC
  & \madeupresult{88.4} & \madeupresult{88.1} & \madeupresult{88.6}
  & \madeupresult{$88.4\pm0.3$} \\
Poisoned ASR
  & \madeupresult{58.7} & \madeupresult{61.2} & \madeupresult{56.9}
  & \madeupresult{$58.9\pm2.2$} \\
PRCL ACC
  & \madeupresult{87.5} & \madeupresult{87.2} & \madeupresult{87.8}
  & \madeupresult{$87.5\pm0.3$} \\
PRCL ASR
  & \madeupresult{18.3} & \madeupresult{20.1} & \madeupresult{16.8}
  & \madeupresult{$18.4\pm1.7$} \\
\bottomrule
\end{tabular}
\end{table}

Variance across seeds is expected to be low ($\sigma\approx 0.3$ for ACC,
$\sigma\approx 1.7$ for ASR), indicating stable training.


% ============================================================================
\section{Theoretical Analysis (Sketch)}
\label{sec:theory}

\subsection{Contamination Model}

We model the training data distribution as a contaminated mixture:
\begin{equation}
  X \sim (1-\rho)\cdot P + \rho\cdot Q
\end{equation}
where $P$ is the clean distribution, $Q$ is adversarial contamination,
and $\rho\in[0,1]$ is the contamination fraction.

\subsection{Bounded Influence}

PRCL's positive reweighting and gradient capping bound the influence of any
single sample $x_i$ on the loss:
\begin{equation}
  \left|\frac{\partial\mathcal{L}_{\text{PRCL}}}{\partial\theta}\bigg|_{x_i}\right|
  \leq \frac{w_{\min}\cdot C}{\sum_j w_{\text{pos}}(j)}
  = O\!\left(\frac{w_{\min}\cdot C}{N}\right)
\end{equation}
where $C$ is the gradient cap value.  This contrasts with standard InfoNCE,
where a single adversarially positioned sample can contribute $O(1/B)$ to the
batch gradient.

\subsection{Target Guarantee (Informal)}

Under assumptions on PCF quality (AUC $\geq 1-\delta$), PRCL's objective
deviation satisfies:
\begin{equation}
  \bigl|\hat{\mathcal{L}}_{\text{PRCL}}(\theta)
  - \mathcal{L}(\theta;P)\bigr|
  \leq O(\rho) + \varepsilon_n
\end{equation}
A rigorous proof is deferred to the final version.


% ============================================================================
\section{Limitations and Discussion}
\label{sec:limitations}

\subsection{Scope of Claims}

PRCL provides robustness improvements under specific contamination and attack
settings.  We do \textbf{not} claim:
\begin{itemize}[nosep]
  \item Universal protection against all possible attacks
  \item Provable robustness guarantees (partial theoretical analysis only)
  \item Effectiveness against sophisticated adaptive attackers
  \item Protection at $\rho>10\%$
\end{itemize}

\subsection{Adaptive Attacks}

An adaptive attacker who knows PRCL's internals could:
\begin{enumerate}[nosep]
  \item Craft \emph{probe-resilient triggers} stable under blur/occlusion
  \item Distribute triggers to produce low individual suspicion scores
  \item Manipulate rolling z-score statistics
\end{enumerate}
The probe registry is extensible---defenders can add custom probes for specific
threat scenarios.

\subsection{Computational Overhead}

PRCL introduces $K$ additional forward passes per batch (one per probe).  With
the default $K=2$ probes, this adds $\sim$40--60\% overhead vs.\ standard
SimCLR---acceptable for security-critical applications but potentially
prohibitive for very large-scale pretraining.

\subsection{Scaling}

Evaluation is limited to CIFAR-scale datasets and ResNet backbones.  Scaling to
ImageNet-1k and Vision Transformers remains future work.


% ============================================================================
\section{Reproducibility}
\label{sec:reproducibility}

\subsection{Environment Setup}

\begin{lstlisting}[language=bash]
git clone https://github.com/aayushakumar/PRCL.git
cd PRCL
pip install -e '.[dev]'
python -m pytest tests/ -x -q   # 99 tests pass
\end{lstlisting}

\subsection{Running Experiments}

\begin{lstlisting}[language=bash]
# Smoke test (< 5 min on T4)
export PRCL_ALLOW_ATTACKS=1
python scripts/train.py ssl.epochs=5 dataset.subset_size=5000 \
    attack=patch_backdoor defense=prcl

# Full main tables
bash scripts/reproduce_main_tables.sh

# Component ablation
bash scripts/ablation.sh

# Poison-ratio sweep
bash scripts/sweep_poison_ratio.sh

# Generate aggregate report
python scripts/report.py --runs-dir ./runs
\end{lstlisting}

\subsection{Run Artifacts}

Every experiment produces structured artifacts in \texttt{runs/}:

\begin{table}[h]
\centering
\caption{Artifacts produced per experimental run.}
\begin{tabular}{@{}ll@{}}
\toprule
\textbf{File} & \textbf{Purpose} \\
\midrule
\texttt{metadata.json} & Git hash, seed, CLI command, timestamp \\
\texttt{hardware.json} & CPU, GPU, CUDA version, memory \\
\texttt{env.txt}       & Full \texttt{pip freeze} \\
\texttt{metrics.json}  & Per-epoch metrics + final evaluation \\
\texttt{run\_card.json} & Structured experiment summary \\
\texttt{best.ckpt}     & Best model checkpoint \\
\texttt{last.ckpt}     & Final epoch checkpoint \\
\texttt{attack\_metadata.json} & Attack config (when enabled) \\
\bottomrule
\end{tabular}
\end{table}

\subsection{Expected Runtime}

\begin{table}[h]
\centering
\caption{Estimated runtime per experiment. \daggernote}
\begin{tabular}{@{}llc@{}}
\toprule
\textbf{Experiment} & \textbf{GPU} & \textbf{Time} \\
\midrule
Single run (200 epochs, CIFAR-10)  & T4   & \madeupresult{$\sim$3.5h} \\
Single run (200 epochs, CIFAR-10)  & A100 & \madeupresult{$\sim$45min} \\
Full main tables (multi-seed)      & A100 & \madeupresult{$\sim$72h} \\
Full ablation                      & A100 & \madeupresult{$\sim$8h} \\
Smoke test (5 ep, 5k subset)       & T4   & \madeupresult{$\sim$3min} \\
\bottomrule
\end{tabular}
\end{table}


% ============================================================================
\section{Project Timeline and Next Steps}
\label{sec:timeline}

\begin{table}[h]
\centering
\caption{Project timeline.}
\begin{tabular}{@{}lll@{}}
\toprule
\textbf{Phase} & \textbf{Deliverable} & \textbf{Status} \\
\midrule
Phase 1 & Project bootstrap, configs, package structure & \checkmark\, Complete \\
Phase 2 & SimCLR SSL core + unit tests & \checkmark\, Complete \\
Phase 3 & PRCL defense pipeline + tests & \checkmark\, Complete \\
Phase 4 & Attack adapters + IntegritySuite & \checkmark\, Complete \\
Phase 5 & Integration tests (99 total) & \checkmark\, Complete \\
Phase 6 & Docs, scripts, notebook, pre-commit & \checkmark\, Complete \\
\midrule
Phase 7 & GPU experiments (main tables) & Pending \\
Phase 8 & Ablation \& sensitivity sweeps & Pending \\
Phase 9 & Final paper with real results & Pending \\
\bottomrule
\end{tabular}
\end{table}

\subsection{Immediate Next Steps}

\begin{enumerate}
  \item \textbf{Run clean baselines} on CIFAR-10/100 and STL-10 with 3 seeds
    each to establish reference ACC values.
  \item \textbf{Run poisoned-undefended experiments} across both attack families
    and all contamination ratios to measure ASR ceilings.
  \item \textbf{Run PRCL-defended experiments} with default hyperparameters to
    populate Tables~\ref{tab:cifar10_patch}--\ref{tab:stl10}.
  \item \textbf{Run ablation matrix} (Table~\ref{tab:ablation}) and sensitivity
    sweeps (Tables~\ref{tab:lambda_pos}--\ref{tab:lambda_neg}).
  \item \textbf{Generate PCF AUC scores} using forensics module to populate
    Table~\ref{tab:pcf_auc}.
  \item \textbf{Generate visualizations:} suspicion histograms, training curves,
    ASR comparison charts.
  \item \textbf{Write final report} replacing all made-up values with real
    experimental results.
\end{enumerate}


% ============================================================================
\section{Ethics and Responsible Use}
\label{sec:ethics}

This project is designed and released as \textbf{defense research}.  All attack
adapters implement well-known, previously published attack
patterns~\cite{gu2017badnets} and do not introduce novel offensive capabilities.

\textbf{Safety mechanisms:}
\begin{itemize}[nosep]
  \item Attack execution is dual-gated: requires both
    \texttt{PRCL\_ALLOW\_ATTACKS=1} environment variable \emph{and}
    \texttt{attack.enabled:~true} in config.
  \item Safe defaults: attacks disabled, no one-click offensive pipeline.
  \item All attack metadata is logged for forensic auditability.
\end{itemize}


% ============================================================================
\section{Conclusion}
\label{sec:conclusion}

We presented PRCL, a defense framework for contrastive self-supervised learning
that operates entirely during pretraining without requiring labels, a clean
validation set, or knowledge of the specific attack.  PRCL combines
Probe-Consistency Forensics for label-free suspicion scoring with robust
contrastive optimization that bounds the influence of suspicious samples.

\textbf{Current status:} The complete codebase (defense, attacks, evaluation,
IntegritySuite, 99 tests, documentation) is implemented and committed.
Illustrative results suggest PRCL can reduce ASR by
\madeupresult{$\sim$40pp} at $\rho=1\%$ while maintaining clean accuracy
within \madeupresult{$\sim$1.3\%} of baseline.  GPU experiments to obtain real
results are the immediate next step.

\textbf{Future work:}
\begin{itemize}[nosep]
  \item Scaling to ImageNet and Vision Transformers
  \item Adaptive attack evaluation and defense hardening
  \item Formal robustness guarantees under the contamination model
  \item Integration with non-contrastive SSL methods (BYOL, VICReg)
  \item Extension to multimodal (vision-language) pretraining
\end{itemize}


% ============================================================================
\bibliographystyle{plainnat}
\begin{thebibliography}{13}

\bibitem[Chen et~al.(2020)]{chen2020simclr}
T.~Chen, S.~Kornblith, M.~Norouzi, and G.~Hinton.
\newblock A simple framework for contrastive learning of visual representations.
\newblock In \emph{ICML}, 2020.

\bibitem[Grill et~al.(2020)]{grill2020byol}
J.-B. Grill et~al.
\newblock Bootstrap your own latent: A new approach to self-supervised learning.
\newblock In \emph{NeurIPS}, 2020.

\bibitem[Gu et~al.(2017)]{gu2017badnets}
T.~Gu, B.~Dolan-Gavitt, and S.~Garg.
\newblock {BadNets}: Identifying vulnerabilities in the machine learning model
  supply chain.
\newblock \emph{arXiv:1708.06733}, 2017.

\bibitem[Hampel et~al.(1986)]{hampel1986robust}
F.~R. Hampel, E.~M. Ronchetti, P.~J. Rousseeuw, and W.~A. Stahel.
\newblock \emph{Robust Statistics: The Approach Based on Influence Functions}.
\newblock Wiley, 1986.

\bibitem[He et~al.(2020)]{he2020moco}
K.~He, H.~Fan, Y.~Wu, S.~Xie, and R.~Girshick.
\newblock Momentum contrast for unsupervised visual representation learning.
\newblock In \emph{CVPR}, 2020.

\bibitem[Jia et~al.(2022)]{jia2022badencoder}
J.~Jia, Y.~Liu, and N.~Z. Gong.
\newblock {BadEncoder}: Backdoor attacks to pre-trained encoders in
  self-supervised learning.
\newblock In \emph{IEEE S\&P}, 2022.

\bibitem[Li et~al.(2021)]{li2021anti}
Y.~Li, X.~Lyu, N.~Koren, L.~Lyu, B.~Li, and X.~Ma.
\newblock Anti-backdoor learning: Training clean models on poisoned data.
\newblock In \emph{NeurIPS}, 2021.

\bibitem[Oord et~al.(2018)]{oord2018cpc}
A.~van~den Oord, Y.~Li, and O.~Vinyals.
\newblock Representation learning with contrastive predictive coding.
\newblock \emph{arXiv:1807.03748}, 2018.

\bibitem[Saha et~al.(2022)]{saha2022backdoor}
A.~Saha, A.~Subramanya, and H.~Pirsiavash.
\newblock Backdoor attacks on self-supervised learning.
\newblock In \emph{CVPR}, 2022.

\bibitem[Tran et~al.(2018)]{tran2018spectral}
B.~Tran, J.~Li, and A.~Madry.
\newblock Spectral signatures in backdoor attacks.
\newblock In \emph{NeurIPS}, 2018.

\bibitem[Tukey(1960)]{tukey1960survey}
J.~W. Tukey.
\newblock A survey of sampling from contaminated distributions.
\newblock \emph{Contributions to Probability and Statistics}, 1960.

\bibitem[Wang et~al.(2019)]{wang2019neural}
B.~Wang et~al.
\newblock Neural cleanse: Identifying and mitigating backdoor attacks in neural
  networks.
\newblock In \emph{IEEE S\&P}, 2019.

\bibitem[Zhang et~al.(2023)]{zhang2023decree}
Z.~Zhang et~al.
\newblock {DECREE}: Detecting backdoor attacks in self-supervised learning via
  contrastive representation entropy.
\newblock \emph{arXiv}, 2023.

\end{thebibliography}


% ============================================================================
\appendix

\section{Repository Structure}
\label{app:repo}

\begin{lstlisting}[basicstyle=\ttfamily\footnotesize]
PRCL/
+-- configs/                          # Hydra YAML configs
|   +-- config.yaml                   # Top-level defaults
|   +-- attack/                       # none, patch_backdoor, blend_backdoor
|   +-- dataset/                      # cifar10, cifar100, stl10
|   +-- defense/                      # none, prcl
|   +-- eval/                         # linear_probe
|   +-- model/                        # resnet18, resnet50
|   +-- ssl/                          # simclr
+-- src/prcl/                         # Main package
|   +-- ssl/                          # SimCLR (backbone, head, loss, loop)
|   +-- defenses/prcl/                # PCF, robust loss, quarantine
|   +-- attacks/                      # Adapters (safety-gated)
|   +-- eval/                         # Linear probe, ASR, forensics
|   +-- datasets/                     # CIFAR, STL-10, poison wrapper
|   +-- integritysuite/               # Run manager, cards, plots
+-- scripts/                          # train.py, sweeps, report
+-- tests/                            # 99 tests (unit + integration + smoke)
+-- notebooks/                        # Colab notebook
+-- docs/                             # Threat model, ethics, reproducibility
+-- pyproject.toml                    # Package config
\end{lstlisting}

\section{PRCL Default Configuration}
\label{app:config}

\begin{lstlisting}[language=Python,basicstyle=\ttfamily\footnotesize]
# configs/defense/prcl.yaml
name: prcl
enabled: true
pcf:
  enabled: true
  stat: probe_alignment
  k_probes: 2
  probe_types: [blur, occlusion]
  normalize: rolling_zscore
  neighbor_k: 10
robust:
  mode: soft_weight
  lambda_pos: 0.5
  lambda_neg: 0.4
  w_min: 0.2
  trim_alpha: 0.0
  grad_cap_enabled: true
  grad_cap_value: 5.0
negatives:
  reweight: true
  false_negative_correction: false
quarantine:
  enabled: false
  threshold: 0.95
  reevaluate_every: 5
  percentile: 0.99
\end{lstlisting}

\section{Hydra CLI Examples}
\label{app:cli}

\begin{lstlisting}[language=bash,basicstyle=\ttfamily\footnotesize]
# Clean baseline
python scripts/train.py defense=none attack=none

# Poisoned, no defense
export PRCL_ALLOW_ATTACKS=1
python scripts/train.py defense=none attack=patch_backdoor \
    attack.poison_ratio=0.01

# PRCL defended
python scripts/train.py defense=prcl attack=patch_backdoor \
    attack.poison_ratio=0.01

# Override parameters
python scripts/train.py defense=prcl \
    defense.robust.lambda_pos=0.3 ssl.epochs=100 seed=123
\end{lstlisting}


\end{document}
